%% Unary normal modal completeness, with admissible valuations explicit.

\section{Lindenbaum algebras and modal completeness}

\begin{annexitem}{annex:p10.modal-algebra.1}{thm:p10-modal-algebra-completeness}{The modal quotient and algebraic completeness}
\begin{proof}
Classical propositional reasoning makes \(\equiv_L\) an
equivalence relation compatible with all Boolean connectives.
For the modal constructor, if
\(\vdash_L\varphi\leftrightarrow\psi\), then both implications
are theorems.  Necessitate each and apply distribution to obtain
\[
 \vdash_L\Box\varphi\to\Box\psi,\qquad
 \vdash_L\Box\psi\to\Box\varphi.
\]
Their conjunction yields
\(\vdash_L\Box\varphi\leftrightarrow\Box\psi\).
The quotient operations are therefore well-defined.
Boolean identities hold because their formula counterparts are
propositional tautologies.

The formula \(\Box\top\) is a theorem by necessitation.
Necessitating the two projections of a conjunction and applying
distribution proves
\(\Box(\varphi\wedge\psi)\to
(\Box\varphi\wedge\Box\psi)\).
For the converse, necessitate
\(\varphi\to(\psi\to\varphi\wedge\psi)\) and apply
distribution twice.
Thus \(\Box\) preserves \(1\) and binary meets in \(F_L\).
Taking Boolean duals shows that \(\Diamond\) preserves
\(0\) and binary joins.
This proves that \(F_L\) is a normal modal algebra.

Classical tautologies also give
\[
 [\chi]_L=1\quad\Longleftrightarrow\quad\chi\in L:
\]
each of \(\chi\) and \(\chi\leftrightarrow\top\) implies
the other by propositional reasoning.
Induction on formulas gives
\[
 \llbracket\chi\rrbracket_\nu^{F_L}=[\chi]_L,
 \qquad \nu(p)=[p]_L.
\]

\paragraph{All valuations in the quotient.}
Fix \(\lambda\in L\) and an arbitrary valuation
\(v:\mathsf{Prop}\to F_L\).
Only finitely many variables occur in \(\lambda\).
For each of those variables \(p\), select a formula \(\rho_p\)
with \([\rho_p]_L=v(p)\).
Define a substitution \(\sigma\) by these finitely many choices,
and put \(\sigma(q)=q\) on all other variables.
Induction on the subformulas of \(\lambda\) gives
\[
 \llbracket\lambda\rrbracket_v^{F_L}=[\sigma(\lambda)]_L.
\]
Uniform-substitution closure puts \(\sigma(\lambda)\) in \(L\),
so this value is \(1\).
This proves \(F_L\in\operatorname{Alg}(L)\).
There was no selection of representatives for a whole infinite
family of equivalence classes.

If \(\varphi\notin L\), its value under \(\nu\) is not \(1\),
so \(F_L\) refutes it.  Conversely, every member of \(L\)
is valid in every member of \(\operatorname{Alg}(L)\) by
definition.  These statements prove all displayed equivalences.

The algebra definition also respects the stated normal-logic rules.
Meet preservation makes \(\Box_A\) monotone, and
\[
 \Box_A(a\to b)\wedge\Box_Aa
       =\Box_A((a\to b)\wedge a)\leq\Box_Ab
\]
proves distribution.
Modus ponens preserves validity, substitution is composition
with formula evaluation, and theorem necessitation preserves
validity because \(\Box_A1=1\).
Thus, if \(L\) is presented by additional axioms, validating
those axioms in a normal modal algebra validates every theorem.

If \(\bot\in L\), propositional explosion gives every formula
in \(L\); hence every two formulas are equivalent and \(F_L\)
is trivial.  Conversely, a trivial quotient has
\([\bot]_L=1\), so \(L\) is inconsistent.
Any algebra validating this logic must have \(0=1\);
the trivial modal algebras do validate every formula.
\end{proof}
\end{annexitem}

\begin{annexitem}{annex:p10.modal-algebra.2}{thm:p10-modal-general-frame-completeness}{The canonical general frame}
\begin{proof}
Assume \(L\) is consistent.
Then \(F_L\) is nontrivial, so the Boolean prime ideal theorem
gives at least one ultrafilter and \(X_L\ne\varnothing\).
Theorem~\ref{thm:p10-jonsson-tarski} says that
\(a\mapsto\widehat a\) is an injective Boolean and modal
homomorphism into the full complex algebra of \((X_L,R_L)\).
Its image is closed under the required operations and equals
the Stone clopens.  It consequently supplies the admissible
algebra of \(\mathfrak G_L\).

For clarity, the two descriptions of \(R_L\) agree.
If the diamond condition holds and \(\Box a\in U\) but
\(a\notin V\), then \(\neg a\in V\) gives
\(\Diamond\neg a\in U\), contradicting
\(\Box a=\neg\Diamond\neg a\).
Conversely, if the box condition holds and \(a\in V\)
but \(\Diamond a\notin U\), then
\(\Box\neg a=\neg\Diamond a\in U\), forcing
\(\neg a\in V\), again a contradiction.

For an arbitrary admissible valuation \(w\), every \(w(p)\)
has a unique preimage \(v(p)\in F_L\) under the injective
Stone map.  This defines a function by uniqueness, without
choosing formula representatives.
Write \(\llbracket\chi\rrbracket_w^{\mathfrak G_L}\) for the set
of worlds satisfying \(\chi\) under \(w\).
Induction on formulas, using preservation of the modal operation,
proves
\[
 \llbracket\chi\rrbracket_w^{\mathfrak G_L}
       =\widehat{\llbracket\chi\rrbracket_v^{F_L}}.
\]
Since \(F_L\) validates all of \(L\) under every valuation,
\(\mathfrak G_L\) validates all of \(L\) under every admissible
valuation.  The finite-representative argument needed for that
algebraic validity was proved in the preceding annex item.
Taking \(v=\nu\) gives the canonical truth lemma.

If \(\varphi\notin L\), then \(\neg[\varphi]_L\ne0\).
Its principal filter is proper; extend it to an ultrafilter \(U\).
The truth lemma refutes \(\varphi\) at \(U\) under \(v_0\).

\paragraph{Agreement with the Part VIII canonical frame.}
For an ultrafilter \(U\), set
\(\Sigma_U=\{\varphi:[\varphi]_L\in U\}\).
It contains the theorems of \(L\), is closed under modus
ponens, and contains exactly one of \(\varphi,\neg\varphi\).
It is consistent: all lines of a derivation from its members
remain in \(\Sigma_U\) by the closure just proved, whereas a
contradiction would put \(0\) in the proper filter \(U\).
Its decision of each formula makes it maximal consistent.

Conversely, a maximal \(L\)-consistent set \(\Sigma\)
contains all \(L\)-theorems and is closed under finite
derivability: otherwise adding a missing consequence would
preserve consistency and contradict maximality.
It cannot contain both \(\varphi\) and \(\neg\varphi\).
If it contained neither, maximality would make both one-formula
extensions inconsistent; their finite contradiction derivations
combine by propositional case analysis into a contradiction
from \(\Sigma\).
It therefore contains exactly one of the two, and propositional
reasoning shows that it contains a conjunction exactly when it
contains both conjuncts.
It respects \(\equiv_L\), so
\(\{[\varphi]_L:\varphi\in\Sigma\}\) is well-defined.
This is a proper filter: meets are preserved, and
\[
 [\varphi]_L\leq[\psi]_L
      \quad\Longleftrightarrow\quad
      \vdash_L\varphi\to\psi
\]
makes it upward closed.
Its decision of complements makes it an ultrafilter.
The two maps are inverse.
Finally the box description of \(R_L\) becomes precisely
\[
 \Sigma_U R_L\Sigma_V
 \quad\Longleftrightarrow\quad
 \{\chi:\Box\chi\in\Sigma_U\}\subseteq\Sigma_V.
\]
This is the Part VIII relation, with the same membership
valuation for variables.
\end{proof}
\end{annexitem}

\begin{annexitem}{annex:p10.modal-algebra.3}{thm:p10-modal-general-frame-strong}{Strong local completeness}
\begin{proof}
With theorem-only necessitation, the local derivability convention
has the finite deduction property
\[
 \Gamma\vdash_L\varphi
 \quad\Longleftrightarrow\quad
 \text{some finite }\Gamma_0\subseteq\Gamma
 \text{ has }\vdash_L(\bigwedge\Gamma_0)\to\varphi.
\]
For the forward direction, collect the finitely many used premises
and apply the propositional deduction argument along the proof.
Every necessitated line has a separate theorem proof, so it is
already a member of \(L\) and does not require necessitation
of a premise-dependent implication.
For the reverse direction, infer the conjunction from its
premises and apply modus ponens.  The empty conjunction is \(\top\).

Theorems of \(L\) hold at every point of a frame in
\(\operatorname{Gen}(L)\), and modus ponens preserves truth
at a fixed point.  The finite deduction property therefore
proves soundness for local consequence.

Suppose \(\Gamma\not\vdash_L\varphi\).
Then \(L\) is consistent.
For every finite \(\Gamma_0\subseteq\Gamma\),
\[
 \left(\bigwedge_{\gamma\in\Gamma_0}[\gamma]_L\right)
                    \wedge\neg[\varphi]_L\ne0.
\]
Otherwise the corresponding implication would belong to \(L\),
contradicting the deduction property.
Thus the Boolean filter generated by
\[
 \{[\gamma]_L:\gamma\in\Gamma\}
                 \cup\{\neg[\varphi]_L\}
\]
is proper: none of its finite generating meets is zero.
The Boolean prime ideal theorem extends it to an ultrafilter \(U\).
The canonical truth lemma makes all of \(\Gamma\) true at \(U\)
under \(v_0\), and makes \(\varphi\) false there.
The preceding theorem places \(\mathfrak G_L\) in
\(\operatorname{Gen}(L)\), proving completeness.

If \(L\) is inconsistent, every formula is a theorem and hence
derivable from every premise set.
No nonempty general frame validates \(\bot\), so
\(\operatorname{Gen}(L)=\varnothing\); semantic consequence
over that empty class is vacuous.
This handles the remaining case without inserting an empty
frame into Part VIII's nonempty-frame convention.

The algebraic theorem above concerns validity of a single formula.
For finitely many local premises its corresponding algebraic
condition is the inequality
\(\bigwedge_{\gamma\in\Gamma_0}\llbracket\gamma\rrbracket_v
\leq\llbracket\varphi\rrbracket_v\), that is, validity of the
implication.  Merely requiring every premise value to equal
\(1\) would instead lose this distinction: for example,
\(v(p)=1\) forces \(v(\Box p)=1\), although \(p\) does not
locally entail \(\Box p\) on all Kripke frames.
A two-world frame with \(wRv\) and \(p\) true only at \(w\)
gives the latter counterexample.
\end{proof}
\end{annexitem}

\begin{annexitem}{annex:p10.modal-algebra.4}{thm:p10-modal-canonicity-completeness}{Passing to ordinary frames}
\begin{proof}
Assume that the consistent logic \(L\) satisfies
\(\mathfrak F_L\models L\) under unrestricted valuations.
If \(\Gamma\not\vdash_L\varphi\), the preceding proof supplies
an ultrafilter \(U\) and the canonical valuation \(v_0\)
making the premises true and the conclusion false.
Every admissible set is a subset of \(X_L\), so this same
valuation is permitted on the ordinary frame \(\mathfrak F_L\).
The recursive satisfaction clauses are unchanged.
By the assumed validity condition
\(\mathfrak F_L\in\operatorname{Fr}(L)\), and this is the
required ordinary-frame counterexample.
Soundness follows as before from theorems of \(L\) and the
finite deduction property.

Now suppose \(\operatorname{Alg}(L)\) is closed under canonical
extensions.  Since \(F_L\in\operatorname{Alg}(L)\), its full
complex extension \(F_L^\sigma\) also belongs to that class.
Ordinary valuations on \(\mathfrak F_L\) are exactly valuations
in the power-set carrier of \(F_L^\sigma\).
Formula induction identifies their truth sets with their
complex-algebra values.
Hence \(F_L^\sigma\models L\) is precisely
\(\mathfrak F_L\models L\), establishing the sufficient condition.

The argument uses only the displayed fixed-language canonical
frame.  It asserts no converse from this one frame's validity
to closure of the whole algebra class under canonical extensions.
The all-variable-set condition in Venema's formulation addresses
that separate equivalence.
\end{proof}
\end{annexitem}
